% Generated draft: requires mathematical and source review.
\documentclass[11pt]{article}
\usepackage[margin=25mm]{geometry}
\usepackage{fontspec}
\setmainfont{Latin Modern Roman}
\usepackage{amsmath,amssymb,mathtools,unicode-math}
\setmathfont{Latin Modern Math}
\usepackage{enumitem,needspace,xurl,hyperref}
\setlength{\parindent}{0pt}
\setlength{\parskip}{5pt}
\setlength{\emergencystretch}{3em}
\newcommand{\N}{\mathbb{N}}
\newcommand{\Z}{\mathbb{Z}}
\newcommand{\Q}{\mathbb{Q}}
\newcommand{\R}{\mathbb{R}}
\newcommand{\C}{\mathbb{C}}
\newcommand{\F}{\mathbb{F}}
\newcommand{\A}{\mathbb{A}}
\newcommand{\PP}{\mathbb{P}}
\newcommand{\E}{\mathbb{E}}
\newcommand{\eps}{\varepsilon}
\newcommand{\dd}{\,\mathrm{d}}
\begin{document}
\section*{UnsolvedMath ID: 30003083}
\textbf{Source problem code:} catalogue entry
\subsection{Definitions and mathematical statement}
Let \(X\) be a smooth projective algebraic surface over \(\C\). An irreducible curve \(C\subset X\) means a reduced, irreducible closed subvariety of dimension one; no smoothness assumption is imposed on \(C\). Write \(C^2=C\cdot C\in\Z\) for its self-intersection number in the intersection theory of divisors on \(X\).

\textbf{Bounded Negativity Conjecture.} For every smooth complex projective surface \(X\), there exists a nonnegative integer \(b(X)\) such that
\[
C^2\geq -b(X)
\qquad\text{for every irreducible curve }C\subset X.
\]
The bound may depend on \(X\), but must be independent of \(C\). Taking \(b(X)\) to be a nonnegative integer is equivalent to allowing an arbitrary finite real bound, since self-intersection numbers are integers.

The supplied record labels this conjecture open. Its dated background additionally reports results for unspecified broad classes of surfaces and bounded-genus families, and failure in positive characteristic. The hypotheses and scope of those background claims are not supplied; they are recorded as source qualifications and are not additional assertions or extensions of the conjecture stated here.
\subsection{Short English statement}
For each fixed smooth complex projective surface, the self-intersection numbers of all its irreducible curves should have a common finite lower bound. The bound may vary with the surface, and singular irreducible curves are included. This is a statement-only draft requiring human mathematical review.
\subsection{Sources}
\begin{itemize}
\item Supplied local UnsolvedMath registry record, permanent ID \(30003083\), exact source problem code \texttt{catalogue entry}; title: \emph{Bounded Negativity Conjecture for Projective Surfaces}. The record belongs to the Oberwolfach Reports Open Problems collection and gives proposed year 2016, without identifying an individual proposer.
\item Supplied report citation: \emph{Mini-Workshop: Arrangements of Subvarieties, and their Applications in Algebraic Geometry} (2016), Oberwolfach Reports, pp.~653--702, DOI 10.4171/owr/2016/14. Supplied literature-source label: OWR report 14/2016. \url{https://doi.org/10.4171/owr/2016/14}.
\item Supplied literature reference: Bauer et al. \url{https://arxiv.org/abs/1109.1881}.
\item Supplied literature reference: Laface--Pokora. \url{https://arxiv.org/abs/1708.09463}.
\item Literature assessment dated \texttt{2026-08-21}: the supplied curation labels the conjecture open in characteristic zero, reports proofs for broad classes of surfaces and bounded-genus families, and reports failure in positive characteristic. This is supplied dated curation, not independently verified here. The record itself requests independent verification of time-sensitive status and cited claims.
\item The supplied statement-status label is \texttt{corrected\_verified}. Its accompanying verification note reports only safe notation, typography, and TeX-delimiter cleanup while preserving hypotheses, quantifiers, and subquestions; this does not constitute independent mathematical review of this draft.
\end{itemize}
\end{document}
